\begin{abstract}
NetworkPolicy binds permissions to labels, and labels to Pods that are
created and replaced continuously. Engines that materialize an entry per
peer identity at each endpoint carry state proportional to endpoints
times identities, regardless of who communicates. kube-network-policies
(\sys{}) places the work elsewhere. nftables steers uncached traffic of
policy-selected Pods to NFQUEUE, a userspace evaluator decides when a
connection is attempted, and an accepted verdict is stored as a conntrack
label so the kernel forwards the rest of the connection. Local Pod
addresses arrive through the runtime's NRI hooks rather than the API
server. We compare \sys{} with Cilium on a 720-worker cluster creating
35,000 Pods in a bidirectional-mesh identity sweep whose sandboxes block
until a permitted TCP connection succeeds, so post-scheduling latency
bounds time to first permitted communication. \sys{} completes the
35,000-identity mesh tier that Cilium's producer excluded on predicted
policy-map capacity (all Pods Running, 421~s P99, two startup-SLO
failures). At 700 and 3,500 identities its P99 is 5.33 and 2.79~s, within
2.4~s of Cilium, at 2.5--3$\times$ lower achieved Pod throughput. At 7
identities with 5,000 Pods per ReplicaSet its P99 is 151~s against
Cilium's 6.8~s, unexplained. Each tier has one trial and no agent metrics.
We show where the eager cross-product goes and what moving it costs; we
do not show an admission-latency advantage.
\end{abstract}

\section{Introduction}
\label{sec:intro}
A Pod is the unit of isolation in Kubernetes and also the unit of change.
Deployments roll, jobs finish, sandboxes are replaced, and each event
changes who may talk to whom. NetworkPolicy expresses permissions through
Pod and namespace labels, ports, and address ranges; the addresses that
carry traffic are assigned later and independently. An enforcement engine
must join policy with current workload metadata on every node hosting a
selected Pod, isolate a new Pod before its containers start, and keep
doing both while packets arrive. A service that lives for hours amortizes
this setup; a sandbox that lives for seconds does not.

The datapath question is usually framed as lookup speed: chains against
sets, iptables against nftables against eBPF. That misses where the cost
has moved. nftables applies incremental updates proportional to what
changed~\cite{winship2025nftables}; Open vSwitch sends misses to userspace
and caches results in the kernel~\cite{pfaff2015ovs}; eBPF maps let
userspace update keyed state without reloading code~\cite{linuxbpfmaps}.
All three make the lookup cheap. None decides how much state must exist
before a Pod can communicate, or on how many nodes. An engine that stores
one entry per peer identity per direction at each endpoint carries state
proportional to identities times local endpoints
(Section~\ref{sec:model}), and each new identity is a write to every
endpoint whose policy selects it. Identity sharing among equally labeled
Pods limits this; label cardinality and broad selectors restore it. What
matters is distinct identities, permission-graph density, and churn, not
Pod count.

\sys{} does not materialize permitted peers. nftables records which local
addresses are subject to policy and sends packets with no cached decision
to NFQUEUE. A userspace evaluator resolves both addresses to Pod metadata,
applies the NetworkPolicy semantics, and returns a verdict. An accept
verdict also sets a conntrack label, and a kernel rule then accepts any
packet whose connection carries the label in established or related
state, so the agent sees a connection roughly once. Denied packets
receive no label. Local Pod addresses reach the agent through the
runtime's NRI \texttt{RunPodSandbox} hook without waiting for the kubelet
to publish them; remote metadata still arrives through API watches.

The research question is whether this trade is favorable for clusters
with many identities and high churn, not whether userspace is faster than
the kernel; per packet it is slower. The trade replaces a cross-product
of identities and endpoints, maintained eagerly on every node, with
per-connection work proportional to communication actually attempted,
plus metadata and conntrack that both designs carry. It also introduces a
bounded queue whose failure policy defaults to open
(Section~\ref{sec:design}).

Our evidence is archived ClusterLoader2 runs on a 720-worker
\texttt{n2-standard-4} cluster creating 35,000 Pods at a configured
500~QPS, comparing Cilium v1.20.1 with Kindnet v1.0.1, which bundles
\sys{}. Every sandbox has a \texttt{wait-for-gateway} init container that
loops until a permitted TCP connection to the gateway Service succeeds
under default deny, so a Pod is Running only after the source node's
egress evaluation and the gateway node's ingress evaluation have both
admitted its first connection. The post-scheduling latency we report
therefore bounds time to first permitted communication, including
propagation of the new Pod's address to the gateway node. The evidence
supports five observations, none of which stands in for another.
\begin{enumerate}
\setlength{\itemsep}{0pt}
\item Cilium's completed identity sweeps show no rising latency curve:
open-gateway P99 after scheduling falls from 5.930~s at 7 identities to
1.990~s at 35,000, hub-and-spoke from 6.370 to 2.033~s, and mean worker
CPU stays within 0.51--0.63 cores in every tier of every family.
\item The 35,000-ReplicaSet tiers halve achieved Phase-1 Pod throughput,
from roughly 280--330 to 170 Pods/s, and raise cluster-scoped Pod LIST P99
from 14--19 to 28.8~s; the bottleneck there is controller fan-out, not
the datapath.
\item Cilium's bidirectional mesh completes at 7, 700, and 3,500
identities (6.767, 2.934, 2.279~s P99). Its 35,000-identity tier was
excluded before execution: two entries per peer identity give 70,000
entries per endpoint against a 65,536-entry policy-map maximum. That is a
predicted capacity exclusion, not an observed failure.
\item Two Cilium admission failures are recorded: v1.18.6 stranded 19 of
35,000 Pods at 700 identities, attributed to a policy-cache race fixed in
v1.20.0~\cite{cilium7515}, and a v1.20.1 hub-and-spoke raw-Pod burst
aborted with 29,306 Running and 5,694 stranded.
\item \sys{}'s mesh sweep, matched to Cilium's generator at the first
three tiers, reaches 151.479, 5.331, and 2.791~s P99. Its unmatched
35,000-identity raw-Pod mesh tier completes with every Pod Running at
421.410~s P99 and two startup-SLO failures. At 700 and 3,500 identities
\sys{} created Pods at 103--124 Pods/s where Cilium created 302--328 on
the same configured QPS; the artifacts do not record why.
\end{enumerate}

We do not claim an admission-latency advantage over Cilium; the evidence
does not show one. \sys{} is slower at every matched tier, by 0.5--2.4~s
at 700 and 3,500 identities and by 22$\times$ at 7. The throughput gap
confounds the comparison in \sys{}'s favor, since a slower generator
produces smaller per-node bursts. Each tier has one trial. Every
agent-level Prometheus query returned no samples because the deployed
Kindnet build exposes no metrics port, so verdict rate, queue occupancy,
drops, and agent CPU are absent. The sweep establishes that \sys{} has no
per-endpoint state bound that excludes the dense tier, and that its
post-scheduling P99 does not rise from 700 to 3,500 identities.

The paper makes three contributions at three levels of evidence: from
source inspection of revision \texttt{9984dcd}, a description of the
enforcement decomposition, its queue-overflow trade-off, and its
unknown-peer semantics (Sections~\ref{sec:design}
and~\ref{sec:implementation}); an analytical model of where eager
permission state goes under deferred evaluation and of the admission
capacity that replaces it (Section~\ref{sec:model}), whose single
service-rate data point is a repository microbenchmark of unrecorded
provenance; and an analysis of the archived benchmark collection across
both engines, including the newly executed \sys{} sweep
(Section~\ref{sec:evaluation}), with the missing measurements enumerated
in Section~\ref{sec:gaps}. Section~\ref{sec:background} states the
problem, Section~\ref{sec:related} places the design among reactive
caches and factored kernel representations, and
Section~\ref{sec:discussion} covers metadata distribution.
